\documentclass[11pt,a4paper]{article}
\usepackage[utf8]{inputenc}
\usepackage[T1]{fontenc}
\usepackage[spanish,es-noshorthands]{babel}
\usepackage[margin=2.2cm]{geometry}
\usepackage{amsmath}
\usepackage{booktabs}
\usepackage{longtable}
\usepackage{caption}
\usepackage{tikz}
\usetikzlibrary{arrows.meta}
\setlength{\tabcolsep}{4.5pt}

\definecolor{baja}{RGB}{42,120,214}
\definecolor{alta}{RGB}{227,73,72}
\definecolor{neutro}{RGB}{137,135,129}
\definecolor{fantasma}{RGB}{200,199,193}
\definecolor{apoyocol}{RGB}{74,58,167}
\definecolor{cargacol}{RGB}{237,161,0}
\definecolor{momcol}{RGB}{0,150,120}
\definecolor{diagcol}{RGB}{214,95,40}


\title{Trípode espacial}

\author{Método de rigidez directa --- viga de Euler-Bernoulli 3D}

\date{28/09/2026}

\begin{document}

\maketitle

\section{Modelo}

La estructura tiene 10 nodos y 10 elementos de viga con 6 grados de libertad por nodo (60 en total, 24 restringidos). Las inercias pueden variar linealmente dentro de cada elemento; la rigidez a flexión se integra de forma exacta con $I(\xi)=I_1(1-\xi)+I_2\,\xi$.

\begin{center}
\begin{minipage}{\linewidth}
\centering
\captionof{table}{Coordenadas de los nodos.}\label{tab:viga3d-nodos}
\begin{tabular}{@{}cccc@{}}
\toprule
Nodo & $x$ [m] & $y$ [m] & $z$ [m] \\
\midrule
0 & $0.000$ & $0.000$ & $0.000$ \\
1 & $2.000$ & $0.000$ & $0.000$ \\
2 & $2.000$ & $4.000$ & $0.000$ \\
3 & $0.000$ & $4.000$ & $0.000$ \\
4 & $0.000$ & $0.000$ & $2.000$ \\
5 & $2.000$ & $0.000$ & $2.000$ \\
6 & $2.000$ & $4.000$ & $2.000$ \\
7 & $0.000$ & $4.000$ & $2.000$ \\
8 & $1.000$ & $0.000$ & $2.000$ \\
9 & $1.000$ & $4.000$ & $2.000$ \\
\bottomrule
\end{tabular}
\end{minipage}
\end{center}

\begin{center}
\begin{minipage}{\linewidth}
\centering
\captionof{table}{Elementos: conectividad y propiedades (inercias en cm$^4$, ``inicio / fin'' si el área o las inercias varían linealmente).}\label{tab:viga3d-elementos}
\begin{tabular}{@{}cccccccc@{}}
\toprule
Elem. & $i$ & $j$ & $L$ [m] & $A$ [cm$^2$] & $I_x$ & $I_y$ & $J_p$ \\
\midrule
0 & 0 & 8 & $2.236$ & $10.00$ & $8.0$ & $8.0$ & $15.9$ \\
1 & 1 & 5 & $2.000$ & $10.00$ & $8.0$ & $8.0$ & $15.9$ \\
2 & 2 & 6 & $2.000$ & $10.00$ & $8.0$ & $8.0$ & $15.9$ \\
3 & 3 & 9 & $2.236$ & $10.00$ & $8.0$ & $8.0$ & $15.9$ \\
4 & 4 & 8 & $1.000$ & $10.00$ & $8.0$ & $8.0$ & $15.9$ \\
5 & 8 & 5 & $1.000$ & $10.00$ & $8.0$ & $8.0$ & $15.9$ \\
6 & 5 & 6 & $4.000$ & $10.00$ & $8.0$ & $8.0$ & $15.9$ \\
7 & 6 & 9 & $1.000$ & $10.00$ & $8.0$ & $8.0$ & $15.9$ \\
8 & 9 & 7 & $1.000$ & $10.00$ & $8.0$ & $8.0$ & $15.9$ \\
9 & 7 & 4 & $4.000$ & $10.00$ & $8.0$ & $8.0$ & $15.9$ \\
\bottomrule
\end{tabular}
\end{minipage}
\end{center}

\begin{figure}[!tbp]
\centering
\begin{tikzpicture}[x={(1.5293cm,-0.4415cm)}, y={(0.8829cm,0.7646cm)}, z={(0.0000cm,1.5293cm)}]
\draw[alta!73!baja, line width=2.2pt] (0.0000,0.0000,0.0000) -- (1.0000,0.0000,2.0000);
\draw[alta!50!baja, line width=2.2pt] (2.0000,0.0000,0.0000) -- (2.0000,0.0000,2.0000);
\draw[alta!50!baja, line width=2.2pt] (2.0000,4.0000,0.0000) -- (2.0000,4.0000,2.0000);
\draw[alta!53!baja, line width=2.2pt] (0.0000,4.0000,0.0000) -- (1.0000,4.0000,2.0000);
\draw[alta!100!baja, line width=2.2pt] (0.0000,0.0000,2.0000) -- (1.0000,0.0000,2.0000);
\draw[alta!32!baja, line width=2.2pt] (1.0000,0.0000,2.0000) -- (2.0000,0.0000,2.0000);
\draw[alta!31!baja, line width=2.2pt] (2.0000,0.0000,2.0000) -- (2.0000,4.0000,2.0000);
\draw[alta!32!baja, line width=2.2pt] (2.0000,4.0000,2.0000) -- (1.0000,4.0000,2.0000);
\draw[alta!100!baja, line width=2.2pt] (1.0000,4.0000,2.0000) -- (0.0000,4.0000,2.0000);
\draw[alta!20!baja, line width=2.2pt] (0.0000,4.0000,2.0000) -- (0.0000,0.0000,2.0000);
\draw[black!80, line width=0.7pt] plot[smooth] coordinates {(0.0000,0.0000,0.0000) (0.0651,-0.0007,0.1237) (0.1347,-0.0026,0.2451) (0.2081,-0.0053,0.3647) (0.2844,-0.0086,0.4828) (0.3628,-0.0123,0.5998) (0.4425,-0.0159,0.7162) (0.5227,-0.0192,0.8323) (0.6026,-0.0220,0.9486) (0.6813,-0.0238,1.0655) (0.7581,-0.0245,1.1833) (0.8320,-0.0237,1.3026) (0.9024,-0.0211,1.4237) (0.9683,-0.0165,1.5469) (1.0290,-0.0095,1.6728) (1.0836,0.0002,1.8018) (1.1313,0.0128,1.9341)};
\draw[black!80, line width=0.7pt] plot[smooth] coordinates {(2.0000,0.0000,0.0000) (2.0014,-0.0009,0.1250) (2.0053,-0.0034,0.2500) (2.0114,-0.0071,0.3750) (2.0193,-0.0116,0.5000) (2.0287,-0.0166,0.6250) (2.0392,-0.0217,0.7500) (2.0506,-0.0266,0.8749) (2.0624,-0.0309,0.9999) (2.0743,-0.0342,1.1249) (2.0859,-0.0362,1.2499) (2.0970,-0.0365,1.3749) (2.1071,-0.0347,1.4999) (2.1159,-0.0306,1.6249) (2.1231,-0.0236,1.7499) (2.1284,-0.0135,1.8749) (2.1313,0.0000,1.9999)};
\draw[black!80, line width=0.7pt] plot[smooth] coordinates {(2.0000,4.0000,0.0000) (2.0014,4.0009,0.1250) (2.0053,4.0034,0.2500) (2.0114,4.0071,0.3750) (2.0193,4.0116,0.5000) (2.0287,4.0166,0.6250) (2.0392,4.0217,0.7500) (2.0506,4.0266,0.8749) (2.0624,4.0309,0.9999) (2.0743,4.0342,1.1249) (2.0859,4.0362,1.2499) (2.0970,4.0365,1.3749) (2.1071,4.0347,1.4999) (2.1159,4.0306,1.6249) (2.1231,4.0236,1.7499) (2.1284,4.0135,1.8749) (2.1313,4.0000,1.9999)};
\draw[black!80, line width=0.7pt] plot[smooth] coordinates {(0.0000,4.0000,0.0000) (0.0651,4.0007,0.1237) (0.1347,4.0026,0.2451) (0.2081,4.0053,0.3647) (0.2844,4.0086,0.4828) (0.3628,4.0123,0.5998) (0.4425,4.0159,0.7162) (0.5227,4.0192,0.8323) (0.6026,4.0220,0.9486) (0.6813,4.0238,1.0655) (0.7581,4.0245,1.1833) (0.8320,4.0237,1.3026) (0.9024,4.0211,1.4237) (0.9683,4.0165,1.5469) (1.0290,4.0095,1.6728) (1.0836,3.9998,1.8018) (1.1313,3.9872,1.9341)};
\draw[black!80, line width=0.7pt] plot[smooth] coordinates {(0.1313,0.0000,1.6005) (0.1938,0.0014,1.6268) (0.2563,0.0027,1.6531) (0.3188,0.0041,1.6792) (0.3813,0.0054,1.7050) (0.4438,0.0067,1.7302) (0.5063,0.0080,1.7548) (0.5688,0.0091,1.7787) (0.6313,0.0102,1.8015) (0.6938,0.0111,1.8234) (0.7563,0.0119,1.8440) (0.8188,0.0126,1.8632) (0.8813,0.0131,1.8810) (0.9438,0.0133,1.8971) (1.0063,0.0134,1.9114) (1.0688,0.0132,1.9238) (1.1313,0.0128,1.9341)};
\draw[black!80, line width=0.7pt] plot[smooth] coordinates {(1.1313,0.0128,1.9341) (1.1938,0.0122,1.9430) (1.2563,0.0115,1.9513) (1.3188,0.0108,1.9589) (1.3813,0.0100,1.9658) (1.4438,0.0093,1.9721) (1.5063,0.0084,1.9777) (1.5688,0.0076,1.9827) (1.6313,0.0067,1.9871) (1.6938,0.0059,1.9909) (1.7563,0.0050,1.9940) (1.8188,0.0041,1.9965) (1.8813,0.0033,1.9984) (1.9438,0.0024,1.9996) (2.0063,0.0016,2.0003) (2.0688,0.0008,2.0004) (2.1313,0.0000,1.9999)};
\draw[black!80, line width=0.7pt] plot[smooth] coordinates {(2.1313,0.0000,1.9999) (2.1341,0.2500,1.9709) (2.1365,0.5000,1.9458) (2.1385,0.7500,1.9246) (2.1402,1.0000,1.9072) (2.1415,1.2500,1.8937) (2.1424,1.5000,1.8840) (2.1430,1.7500,1.8782) (2.1432,2.0000,1.8763) (2.1430,2.2500,1.8782) (2.1424,2.5000,1.8840) (2.1415,2.7500,1.8937) (2.1402,3.0000,1.9072) (2.1385,3.2500,1.9246) (2.1365,3.5000,1.9458) (2.1341,3.7500,1.9709) (2.1313,4.0000,1.9999)};
\draw[black!80, line width=0.7pt] plot[smooth] coordinates {(2.1313,4.0000,1.9999) (2.0688,3.9992,2.0004) (2.0063,3.9984,2.0003) (1.9438,3.9976,1.9996) (1.8813,3.9967,1.9984) (1.8188,3.9959,1.9965) (1.7563,3.9950,1.9940) (1.6938,3.9941,1.9909) (1.6313,3.9933,1.9871) (1.5688,3.9924,1.9827) (1.5063,3.9916,1.9777) (1.4438,3.9907,1.9721) (1.3813,3.9900,1.9658) (1.3188,3.9892,1.9589) (1.2563,3.9885,1.9513) (1.1938,3.9878,1.9430) (1.1313,3.9872,1.9341)};
\draw[black!80, line width=0.7pt] plot[smooth] coordinates {(1.1313,3.9872,1.9341) (1.0688,3.9868,1.9238) (1.0063,3.9866,1.9114) (0.9438,3.9867,1.8971) (0.8813,3.9869,1.8810) (0.8188,3.9874,1.8632) (0.7563,3.9881,1.8440) (0.6938,3.9889,1.8234) (0.6313,3.9898,1.8015) (0.5688,3.9909,1.7787) (0.5063,3.9920,1.7548) (0.4438,3.9933,1.7302) (0.3813,3.9946,1.7050) (0.3188,3.9959,1.6792) (0.2563,3.9973,1.6531) (0.1938,3.9986,1.6268) (0.1313,4.0000,1.6005)};
\draw[black!80, line width=0.7pt] plot[smooth] coordinates {(0.1313,4.0000,1.6005) (0.1263,3.7500,1.5259) (0.1221,3.5000,1.4612) (0.1184,3.2500,1.4065) (0.1155,3.0000,1.3618) (0.1131,2.7500,1.3270) (0.1115,2.5000,1.3021) (0.1105,2.2500,1.2872) (0.1102,2.0000,1.2822) (0.1105,1.7500,1.2872) (0.1115,1.5000,1.3021) (0.1131,1.2500,1.3270) (0.1155,1.0000,1.3618) (0.1184,0.7500,1.4065) (0.1221,0.5000,1.4612) (0.1263,0.2500,1.5259) (0.1313,0.0000,1.6005)};
\draw[apoyocol, fill=apoyocol!30, line width=0.6pt] (0.0000,0.0000,0.0000) ++(-0.16cm,-0.16cm) rectangle ++(0.32cm,0.32cm);
\draw[apoyocol, fill=apoyocol!30, line width=0.6pt] (2.0000,0.0000,0.0000) ++(-0.16cm,-0.16cm) rectangle ++(0.32cm,0.32cm);
\draw[apoyocol, fill=apoyocol!30, line width=0.6pt] (2.0000,4.0000,0.0000) ++(-0.16cm,-0.16cm) rectangle ++(0.32cm,0.32cm);
\draw[apoyocol, fill=apoyocol!30, line width=0.6pt] (0.0000,4.0000,0.0000) ++(-0.16cm,-0.16cm) rectangle ++(0.32cm,0.32cm);
\draw[cargacol, line width=1pt, -{Latex}] (0.0000,-0.0000,0.4750) -- (0.0000,0.0000,0.0000);
\node[cargacol, font=\scriptsize, fill=white, inner sep=1pt] at (0.0000,-0.0000,0.4750) {0.0860687~kN};
\draw[cargacol, line width=1pt, -{Latex}] (2.0000,0.0000,0.4629) -- (2.0000,0.0000,0.0000);
\node[cargacol, font=\scriptsize, fill=white, inner sep=1pt] at (2.0000,0.0000,0.4629) {0.0769822~kN};
\draw[cargacol, line width=1pt, -{Latex}] (2.0000,4.0000,0.4629) -- (2.0000,4.0000,0.0000);
\node[cargacol, font=\scriptsize, fill=white, inner sep=1pt] at (2.0000,4.0000,0.4629) {0.0769822~kN};
\draw[cargacol, line width=1pt, -{Latex}] (0.0000,4.0000,0.4750) -- (0.0000,4.0000,0.0000);
\node[cargacol, font=\scriptsize, fill=white, inner sep=1pt] at (0.0000,4.0000,0.4750) {0.0860687~kN};
\draw[cargacol, line width=1pt, -{Latex}] (0.0000,0.0000,2.6171) -- (0.0000,0.0000,2.0000);
\node[cargacol, font=\scriptsize, fill=white, inner sep=1pt] at (0.0000,0.0000,2.6171) {0.192456~kN};
\draw[cargacol, line width=1pt, -{Latex}] (2.0000,0.0000,2.7200) -- (2.0000,0.0000,2.0000);
\node[cargacol, font=\scriptsize, fill=white, inner sep=1pt] at (2.0000,0.0000,2.7200) {0.269438~kN};
\draw[cargacol, line width=1pt, -{Latex}] (2.0000,4.0000,2.7200) -- (2.0000,4.0000,2.0000);
\node[cargacol, font=\scriptsize, fill=white, inner sep=1pt] at (2.0000,4.0000,2.7200) {0.269438~kN};
\draw[cargacol, line width=1pt, -{Latex}] (0.0000,4.0000,2.6171) -- (0.0000,4.0000,2.0000);
\node[cargacol, font=\scriptsize, fill=white, inner sep=1pt] at (0.0000,4.0000,2.6171) {0.192456~kN};
\draw[cargacol, line width=1pt, -{Latex}] (1.0000,-0.0000,2.5779) -- (1.0000,0.0000,2.0000);
\node[cargacol, font=\scriptsize, fill=white, inner sep=1pt] at (1.0000,-0.0000,2.5779) {0.163051~kN};
\draw[cargacol, line width=1pt, -{Latex}] (1.0000,4.0000,2.5779) -- (1.0000,4.0000,2.0000);
\node[cargacol, font=\scriptsize, fill=white, inner sep=1pt] at (1.0000,4.0000,2.5779) {0.163051~kN};
\draw[momcol, line width=1pt, -{Latex}{Latex}] (0.0000,0.0000,0.0000) -- (0.0000,0.4102,-0.0000);
\node[momcol, font=\scriptsize, fill=white, inner sep=1pt] at (0.0000,0.4102,-0.0000) {0.0143448~kN\,m};
\draw[momcol, line width=1pt, -{Latex}{Latex}] (0.0000,4.0000,0.0000) -- (0.0000,4.4102,-0.0000);
\node[momcol, font=\scriptsize, fill=white, inner sep=1pt] at (0.0000,4.4102,-0.0000) {0.0143448~kN\,m};
\draw[momcol, line width=1pt, -{Latex}{Latex}] (0.0000,0.0000,2.0000) -- (-0.7186,0.0449,2.0000);
\node[momcol, font=\scriptsize, fill=white, inner sep=1pt] at (-0.7186,0.0449,2.0000) {0.102843~kN\,m};
\draw[momcol, line width=1pt, -{Latex}{Latex}] (2.0000,0.0000,2.0000) -- (1.2814,-0.0449,2.0000);
\node[momcol, font=\scriptsize, fill=white, inner sep=1pt] at (1.2814,-0.0449,2.0000) {0.102843~kN\,m};
\draw[momcol, line width=1pt, -{Latex}{Latex}] (2.0000,4.0000,2.0000) -- (2.7186,3.9551,2.0000);
\node[momcol, font=\scriptsize, fill=white, inner sep=1pt] at (2.7186,3.9551,2.0000) {0.102843~kN\,m};
\draw[momcol, line width=1pt, -{Latex}{Latex}] (0.0000,4.0000,2.0000) -- (0.7186,4.0449,2.0000);
\node[momcol, font=\scriptsize, fill=white, inner sep=1pt] at (0.7186,4.0449,2.0000) {0.102843~kN\,m};
\draw[momcol, line width=1pt, -{Latex}{Latex}] (1.0000,0.0000,2.0000) -- (1.0000,-0.4102,2.0000);
\node[momcol, font=\scriptsize, fill=white, inner sep=1pt] at (1.0000,-0.4102,2.0000) {0.0143448~kN\,m};
\draw[momcol, line width=1pt, -{Latex}{Latex}] (1.0000,4.0000,2.0000) -- (1.0000,3.5898,2.0000);
\node[momcol, font=\scriptsize, fill=white, inner sep=1pt] at (1.0000,3.5898,2.0000) {0.0143448~kN\,m};
\node[draw, circle, font=\tiny, fill=white, inner sep=0.8pt] at (0.5000,0.0000,1.0000) {0};
\node[draw, circle, font=\tiny, fill=white, inner sep=0.8pt] at (2.0000,0.0000,1.0000) {1};
\node[draw, circle, font=\tiny, fill=white, inner sep=0.8pt] at (2.0000,4.0000,1.0000) {2};
\node[draw, circle, font=\tiny, fill=white, inner sep=0.8pt] at (0.5000,4.0000,1.0000) {3};
\node[draw, circle, font=\tiny, fill=white, inner sep=0.8pt] at (0.5000,0.0000,2.0000) {4};
\node[draw, circle, font=\tiny, fill=white, inner sep=0.8pt] at (1.5000,0.0000,2.0000) {5};
\node[draw, circle, font=\tiny, fill=white, inner sep=0.8pt] at (2.0000,2.0000,2.0000) {6};
\node[draw, circle, font=\tiny, fill=white, inner sep=0.8pt] at (1.5000,4.0000,2.0000) {7};
\node[draw, circle, font=\tiny, fill=white, inner sep=0.8pt] at (0.5000,4.0000,2.0000) {8};
\node[draw, circle, font=\tiny, fill=white, inner sep=0.8pt] at (0.0000,2.0000,2.0000) {9};
\fill (0.0000,0.0000,0.0000) circle (1.4pt);
\node[font=\scriptsize, above left=0pt] at (0.0000,0.0000,0.0000) {0};
\fill (2.0000,0.0000,0.0000) circle (1.4pt);
\node[font=\scriptsize, above left=0pt] at (2.0000,0.0000,0.0000) {1};
\fill (2.0000,4.0000,0.0000) circle (1.4pt);
\node[font=\scriptsize, above left=0pt] at (2.0000,4.0000,0.0000) {2};
\fill (0.0000,4.0000,0.0000) circle (1.4pt);
\node[font=\scriptsize, above left=0pt] at (0.0000,4.0000,0.0000) {3};
\fill (0.0000,0.0000,2.0000) circle (1.4pt);
\node[font=\scriptsize, above left=0pt] at (0.0000,0.0000,2.0000) {4};
\fill (2.0000,0.0000,2.0000) circle (1.4pt);
\node[font=\scriptsize, above left=0pt] at (2.0000,0.0000,2.0000) {5};
\fill (2.0000,4.0000,2.0000) circle (1.4pt);
\node[font=\scriptsize, above left=0pt] at (2.0000,4.0000,2.0000) {6};
\fill (0.0000,4.0000,2.0000) circle (1.4pt);
\node[font=\scriptsize, above left=0pt] at (0.0000,4.0000,2.0000) {7};
\fill (1.0000,0.0000,2.0000) circle (1.4pt);
\node[font=\scriptsize, above left=0pt] at (1.0000,0.0000,2.0000) {8};
\fill (1.0000,4.0000,2.0000) circle (1.4pt);
\node[font=\scriptsize, above left=0pt] at (1.0000,4.0000,2.0000) {9};
\shade[left color=baja, right color=alta] (1.400cm,-1.983cm) rectangle ++(4cm,0.22cm);
\node[font=\scriptsize, anchor=east] at (1.300cm,-1.873cm) {0};
\node[font=\scriptsize, anchor=west] at (5.500cm,-1.873cm) {61.8 MPa};
\node[font=\scriptsize, anchor=north] at (3.400cm,-2.033cm) {$\sigma_{\mathrm{VM}}$ máxima por elemento};
\begin{scope}[shift={(-2.100cm,-1.983cm)}, -{Latex[length=3pt]}, font=\scriptsize]
\draw (0,0,0) -- (0.4530,0.0000,0.0000) node[pos=1.3] {$x$};
\draw (0,0,0) -- (0.0000,0.4530,0.0000) node[pos=1.3] {$y$};
\draw (0,0,0) -- (0.0000,0.0000,0.4530) node[pos=1.3] {$z$};
\end{scope}
\end{tikzpicture}
\caption{Estructura, vista isométrica. Elementos coloreados según $\sigma_{\mathrm{VM}}$; cuadrados: empotramientos, triángulos: apoyos parciales; flechas naranjas: fuerzas, flechas dobles verdes: momentos. En negro fino, la deformada amplificada 50 veces.}
\label{fig:viga3d-estructura}
\end{figure}

\section{Desplazamientos}

El mayor desplazamiento ocurre en el nodo \textbf{4}, con $|u|=8.4111$~mm. La deformada de la figura~\ref{fig:viga3d-estructura} se interpola con las mismas funciones de Hermite de la formulación, así que muestra la curvatura real de cada elemento.

\begin{center}
\begin{minipage}{\linewidth}
\centering
\captionof{table}{Desplazamientos y giros nodales (ejes globales).}\label{tab:viga3d-desplazamientos}
\begin{tabular}{@{}ccccccc@{}}
\toprule
Nodo & $u_x$ [mm] & $u_y$ [mm] & $u_z$ [mm] & $\theta_x$ [mrad] & $\theta_y$ [mrad] & $\theta_z$ [mrad] \\
\midrule
0 & $0.0000$ & $0.0000$ & $0.0000$ & $0.0000$ & $0.0000$ & $0.0000$ \\
1 & $0.0000$ & $0.0000$ & $0.0000$ & $0.0000$ & $0.0000$ & $0.0000$ \\
2 & $0.0000$ & $0.0000$ & $0.0000$ & $0.0000$ & $0.0000$ & $0.0000$ \\
3 & $0.0000$ & $0.0000$ & $0.0000$ & $0.0000$ & $0.0000$ & $0.0000$ \\
4 & $2.6260$ & $0.0003$ & $-7.9907$ & $-6.3651$ & $-8.4344$ & $0.4226$ \\
5 & $2.6257$ & $0.0007$ & $-0.0026$ & $-2.4715$ & $0.2617$ & $-0.2374$ \\
6 & $2.6257$ & $-0.0007$ & $-0.0026$ & $2.4715$ & $0.2617$ & $0.2374$ \\
7 & $2.6260$ & $-0.0003$ & $-7.9907$ & $6.3651$ & $-8.4344$ & $-0.4226$ \\
8 & $2.6260$ & $0.2555$ & $-1.3173$ & $-2.3626$ & $-2.9557$ & $-0.1849$ \\
9 & $2.6260$ & $-0.2555$ & $-1.3173$ & $2.3626$ & $-2.9557$ & $0.1849$ \\
\bottomrule
\end{tabular}
\end{minipage}
\end{center}

\section{Esfuerzos internos}

\begin{center}
\begin{minipage}{\linewidth}
\centering
\captionof{table}{Esfuerzos internos en los extremos de cada elemento (ejes locales; $N>0$ tracción).}\label{tab:viga3d-internas}
\begin{tabular}{@{}cccccccc@{}}
\toprule
Elem. & Nodo & $N$ [kN] & $V_x$ [kN] & $V_y$ [kN] & $T$ [kN\,m] & $M_x$ [kN\,m] & $M_y$ [kN\,m] \\
\midrule
0 & 0 & $-0.427$ & $0.121$ & $-0.093$ & $-0.007$ & $0.070$ & $0.095$ \\
 & 8 & $-0.273$ & $0.051$ & $-0.061$ & $-0.007$ & $-0.103$ & $-0.097$ \\
1 & 1 & $-0.340$ & $0.058$ & $-0.061$ & $-0.001$ & $0.040$ & $0.060$ \\
 & 5 & $-0.186$ & $0.058$ & $-0.061$ & $-0.001$ & $-0.081$ & $-0.056$ \\
2 & 2 & $-0.340$ & $0.058$ & $0.061$ & $0.001$ & $-0.040$ & $0.060$ \\
 & 6 & $-0.186$ & $0.058$ & $0.061$ & $0.001$ & $0.081$ & $-0.056$ \\
3 & 3 & $-0.427$ & $0.150$ & $-0.028$ & $0.007$ & $0.025$ & $0.116$ \\
 & 9 & $-0.273$ & $0.080$ & $0.003$ & $0.007$ & $-0.003$ & $-0.141$ \\
4 & 4 & $0.000$ & $-0.090$ & $0.128$ & $0.050$ & $-0.002$ & $0.002$ \\
 & 8 & $0.000$ & $-0.144$ & $0.182$ & $0.050$ & $0.153$ & $0.120$ \\
5 & 8 & $-0.058$ & $0.026$ & $-0.038$ & $-0.001$ & $0.048$ & $0.040$ \\
 & 5 & $-0.058$ & $-0.029$ & $0.017$ & $-0.001$ & $0.037$ & $0.042$ \\
6 & 5 & $-0.069$ & $0.000$ & $-0.154$ & $0.000$ & $0.082$ & $0.002$ \\
 & 6 & $-0.069$ & $0.000$ & $0.154$ & $0.000$ & $0.082$ & $0.002$ \\
7 & 6 & $-0.058$ & $-0.017$ & $-0.029$ & $0.001$ & $0.042$ & $-0.037$ \\
 & 9 & $-0.058$ & $0.038$ & $0.026$ & $0.001$ & $0.040$ & $-0.048$ \\
8 & 9 & $0.000$ & $-0.182$ & $-0.144$ & $-0.050$ & $0.120$ & $-0.153$ \\
 & 7 & $0.000$ & $-0.128$ & $-0.090$ & $-0.050$ & $0.002$ & $0.002$ \\
9 & 7 & $-0.027$ & $0.000$ & $-0.154$ & $0.000$ & $0.050$ & $0.003$ \\
 & 4 & $-0.027$ & $0.000$ & $0.154$ & $0.000$ & $0.050$ & $0.003$ \\
\bottomrule
\end{tabular}
\end{minipage}
\end{center}

\begin{figure}[!tbp]
\centering
\begin{tikzpicture}[x={(1.5256cm,-0.4404cm)}, y={(0.8808cm,0.7628cm)}, z={(0.0000cm,1.5256cm)}]
\filldraw[diagcol, fill=diagcol!20, line width=0.6pt] (0.0000,0.0000,0.0000) -- (-0.4899,-0.2449,0.2449) -- (0.6869,-0.1566,2.1566) -- (1.0000,0.0000,2.0000) -- cycle;
\draw[black, line width=1.2pt] (0.0000,0.0000,0.0000) -- (1.0000,0.0000,2.0000);
\node[font=\tiny, fill=white, inner sep=0.8pt] at (-0.4899,-0.2449,0.2449) {$-0.43$};
\node[font=\tiny, fill=white, inner sep=0.8pt] at (0.6869,-0.1566,2.1566) {$-0.27$};
\filldraw[diagcol, fill=diagcol!20, line width=0.6pt] (2.0000,0.0000,0.0000) -- (1.5217,0.0000,0.0000) -- (1.7383,0.0000,2.0000) -- (2.0000,0.0000,2.0000) -- cycle;
\draw[black, line width=1.2pt] (2.0000,0.0000,0.0000) -- (2.0000,0.0000,2.0000);
\node[font=\tiny, fill=white, inner sep=0.8pt] at (1.5217,0.0000,0.0000) {$-0.34$};
\node[font=\tiny, fill=white, inner sep=0.8pt] at (1.7383,0.0000,2.0000) {$-0.19$};
\filldraw[diagcol, fill=diagcol!20, line width=0.6pt] (2.0000,4.0000,0.0000) -- (1.5217,4.0000,0.0000) -- (1.7383,4.0000,2.0000) -- (2.0000,4.0000,2.0000) -- cycle;
\draw[black, line width=1.2pt] (2.0000,4.0000,0.0000) -- (2.0000,4.0000,2.0000);
\node[font=\tiny, fill=white, inner sep=0.8pt] at (1.5217,4.0000,0.0000) {$-0.34$};
\node[font=\tiny, fill=white, inner sep=0.8pt] at (1.7383,4.0000,2.0000) {$-0.19$};
\filldraw[diagcol, fill=diagcol!20, line width=0.6pt] (0.0000,4.0000,0.0000) -- (-0.4899,3.7551,0.2449) -- (0.6869,3.8434,2.1566) -- (1.0000,4.0000,2.0000) -- cycle;
\draw[black, line width=1.2pt] (0.0000,4.0000,0.0000) -- (1.0000,4.0000,2.0000);
\node[font=\tiny, fill=white, inner sep=0.8pt] at (-0.4899,3.7551,0.2449) {$-0.43$};
\node[font=\tiny, fill=white, inner sep=0.8pt] at (0.6869,3.8434,2.1566) {$-0.27$};
\filldraw[diagcol, fill=diagcol!20, line width=0.6pt] (0.0000,0.0000,2.0000) -- (0.0000,0.0000,2.0000) -- (1.0000,0.0000,2.0000) -- (1.0000,0.0000,2.0000) -- cycle;
\draw[black, line width=1.2pt] (0.0000,0.0000,2.0000) -- (1.0000,0.0000,2.0000);
\filldraw[diagcol, fill=diagcol!20, line width=0.6pt] (1.0000,0.0000,2.0000) -- (1.0000,-0.0578,2.0578) -- (2.0000,-0.0578,2.0578) -- (2.0000,0.0000,2.0000) -- cycle;
\draw[black, line width=1.2pt] (1.0000,0.0000,2.0000) -- (2.0000,0.0000,2.0000);
\node[font=\tiny, fill=white, inner sep=0.8pt] at (1.0000,-0.0578,2.0578) {$-0.06$};
\node[font=\tiny, fill=white, inner sep=0.8pt] at (2.0000,-0.0578,2.0578) {$-0.06$};
\filldraw[diagcol, fill=diagcol!20, line width=0.6pt] (2.0000,0.0000,2.0000) -- (2.0975,0.0000,2.0000) -- (2.0975,4.0000,2.0000) -- (2.0000,4.0000,2.0000) -- cycle;
\draw[black, line width=1.2pt] (2.0000,0.0000,2.0000) -- (2.0000,4.0000,2.0000);
\node[font=\tiny, fill=white, inner sep=0.8pt] at (2.0975,0.0000,2.0000) {$-0.07$};
\node[font=\tiny, fill=white, inner sep=0.8pt] at (2.0975,4.0000,2.0000) {$-0.07$};
\filldraw[diagcol, fill=diagcol!20, line width=0.6pt] (2.0000,4.0000,2.0000) -- (2.0000,4.0578,1.9422) -- (1.0000,4.0578,1.9422) -- (1.0000,4.0000,2.0000) -- cycle;
\draw[black, line width=1.2pt] (2.0000,4.0000,2.0000) -- (1.0000,4.0000,2.0000);
\node[font=\tiny, fill=white, inner sep=0.8pt] at (2.0000,4.0578,1.9422) {$-0.06$};
\node[font=\tiny, fill=white, inner sep=0.8pt] at (1.0000,4.0578,1.9422) {$-0.06$};
\filldraw[diagcol, fill=diagcol!20, line width=0.6pt] (1.0000,4.0000,2.0000) -- (1.0000,4.0000,2.0000) -- (0.0000,4.0000,2.0000) -- (0.0000,4.0000,2.0000) -- cycle;
\draw[black, line width=1.2pt] (1.0000,4.0000,2.0000) -- (0.0000,4.0000,2.0000);
\filldraw[diagcol, fill=diagcol!20, line width=0.6pt] (0.0000,4.0000,2.0000) -- (-0.0378,4.0000,2.0000) -- (-0.0378,0.0000,2.0000) -- (0.0000,0.0000,2.0000) -- cycle;
\draw[black, line width=1.2pt] (0.0000,4.0000,2.0000) -- (0.0000,0.0000,2.0000);
\node[font=\tiny, fill=white, inner sep=0.8pt] at (-0.0378,4.0000,2.0000) {$-0.03$};
\node[font=\tiny, fill=white, inner sep=0.8pt] at (-0.0378,0.0000,2.0000) {$-0.03$};
\fill (0.0000,0.0000,0.0000) circle (1.3pt);
\node[font=\scriptsize, above left=0pt] at (0.0000,0.0000,0.0000) {0};
\fill (2.0000,0.0000,0.0000) circle (1.3pt);
\node[font=\scriptsize, above left=0pt] at (2.0000,0.0000,0.0000) {1};
\fill (2.0000,4.0000,0.0000) circle (1.3pt);
\node[font=\scriptsize, above left=0pt] at (2.0000,4.0000,0.0000) {2};
\fill (0.0000,4.0000,0.0000) circle (1.3pt);
\node[font=\scriptsize, above left=0pt] at (0.0000,4.0000,0.0000) {3};
\fill (0.0000,0.0000,2.0000) circle (1.3pt);
\node[font=\scriptsize, above left=0pt] at (0.0000,0.0000,2.0000) {4};
\fill (2.0000,0.0000,2.0000) circle (1.3pt);
\node[font=\scriptsize, above left=0pt] at (2.0000,0.0000,2.0000) {5};
\fill (2.0000,4.0000,2.0000) circle (1.3pt);
\node[font=\scriptsize, above left=0pt] at (2.0000,4.0000,2.0000) {6};
\fill (0.0000,4.0000,2.0000) circle (1.3pt);
\node[font=\scriptsize, above left=0pt] at (0.0000,4.0000,2.0000) {7};
\fill (1.0000,0.0000,2.0000) circle (1.3pt);
\node[font=\scriptsize, above left=0pt] at (1.0000,0.0000,2.0000) {8};
\fill (1.0000,4.0000,2.0000) circle (1.3pt);
\node[font=\scriptsize, above left=0pt] at (1.0000,4.0000,2.0000) {9};
\end{tikzpicture}
\caption{Diagrama de $N$ [kN], dibujado hacia $x_L$ positivo para valores positivos.}
\label{fig:viga3d-diag-N}
\end{figure}

\begin{figure}[!tbp]
\centering
\begin{tikzpicture}[x={(1.4622cm,-0.4221cm)}, y={(0.8442cm,0.7311cm)}, z={(0.0000cm,1.4622cm)}]
\filldraw[diagcol, fill=diagcol!20, line width=0.6pt] (0.0000,0.0000,0.0000) -- (-0.1001,0.2503,0.0501) -- (1.1477,-0.3693,1.9261) -- (1.0000,0.0000,2.0000) -- cycle;
\draw[black, line width=1.2pt] (0.0000,0.0000,0.0000) -- (1.0000,0.0000,2.0000);
\node[font=\tiny, fill=white, inner sep=0.8pt] at (-0.1001,0.2503,0.0501) {$0.07$};
\node[font=\tiny, fill=white, inner sep=0.8pt] at (1.1477,-0.3693,1.9261) {$-0.10$};
\filldraw[diagcol, fill=diagcol!20, line width=0.6pt] (2.0000,0.0000,0.0000) -- (2.0000,0.1592,0.0000) -- (2.0000,-0.3185,2.0000) -- (2.0000,0.0000,2.0000) -- cycle;
\draw[black, line width=1.2pt] (2.0000,0.0000,0.0000) -- (2.0000,0.0000,2.0000);
\node[font=\tiny, fill=white, inner sep=0.8pt] at (2.0000,0.1592,0.0000) {$0.04$};
\node[font=\tiny, fill=white, inner sep=0.8pt] at (2.0000,-0.3185,2.0000) {$-0.08$};
\filldraw[diagcol, fill=diagcol!20, line width=0.6pt] (2.0000,4.0000,0.0000) -- (2.0000,3.8408,0.0000) -- (2.0000,4.3185,2.0000) -- (2.0000,4.0000,2.0000) -- cycle;
\draw[black, line width=1.2pt] (2.0000,4.0000,0.0000) -- (2.0000,4.0000,2.0000);
\node[font=\tiny, fill=white, inner sep=0.8pt] at (2.0000,3.8408,0.0000) {$-0.04$};
\node[font=\tiny, fill=white, inner sep=0.8pt] at (2.0000,4.3185,2.0000) {$0.08$};
\filldraw[diagcol, fill=diagcol!20, line width=0.6pt] (0.0000,4.0000,0.0000) -- (-0.0352,4.0881,0.0176) -- (1.0048,3.9880,1.9976) -- (1.0000,4.0000,2.0000) -- cycle;
\draw[black, line width=1.2pt] (0.0000,4.0000,0.0000) -- (1.0000,4.0000,2.0000);
\node[font=\tiny, fill=white, inner sep=0.8pt] at (-0.0352,4.0881,0.0176) {$0.02$};
\node[font=\tiny, fill=white, inner sep=0.8pt] at (1.0048,3.9880,1.9976) {$-0.00$};
\filldraw[diagcol, fill=diagcol!20, line width=0.6pt] (0.0000,0.0000,2.0000) -- (0.0000,-0.0068,1.9932) -- (1.0000,0.4243,2.4243) -- (1.0000,0.0000,2.0000) -- cycle;
\draw[black, line width=1.2pt] (0.0000,0.0000,2.0000) -- (1.0000,0.0000,2.0000);
\node[font=\tiny, fill=white, inner sep=0.8pt] at (0.0000,-0.0068,1.9932) {$-0.00$};
\node[font=\tiny, fill=white, inner sep=0.8pt] at (1.0000,0.4243,2.4243) {$0.15$};
\filldraw[diagcol, fill=diagcol!20, line width=0.6pt] (1.0000,0.0000,2.0000) -- (1.0000,0.1327,2.1327) -- (2.0000,0.1033,2.1033) -- (2.0000,0.0000,2.0000) -- cycle;
\draw[black, line width=1.2pt] (1.0000,0.0000,2.0000) -- (2.0000,0.0000,2.0000);
\node[font=\tiny, fill=white, inner sep=0.8pt] at (1.0000,0.1327,2.1327) {$0.05$};
\node[font=\tiny, fill=white, inner sep=0.8pt] at (2.0000,0.1033,2.1033) {$0.04$};
\filldraw[diagcol, fill=diagcol!20, line width=0.6pt] (2.0000,0.0000,2.0000) -- (2.0000,0.0000,2.3239) -- (2.0000,4.0000,2.3239) -- (2.0000,4.0000,2.0000) -- cycle;
\draw[black, line width=1.2pt] (2.0000,0.0000,2.0000) -- (2.0000,4.0000,2.0000);
\node[font=\tiny, fill=white, inner sep=0.8pt] at (2.0000,0.0000,2.3239) {$0.08$};
\node[font=\tiny, fill=white, inner sep=0.8pt] at (2.0000,4.0000,2.3239) {$0.08$};
\filldraw[diagcol, fill=diagcol!20, line width=0.6pt] (2.0000,4.0000,2.0000) -- (2.0000,4.1168,2.1168) -- (1.0000,4.1124,2.1124) -- (1.0000,4.0000,2.0000) -- cycle;
\draw[black, line width=1.2pt] (2.0000,4.0000,2.0000) -- (1.0000,4.0000,2.0000);
\node[font=\tiny, fill=white, inner sep=0.8pt] at (2.0000,4.1168,2.1168) {$0.04$};
\node[font=\tiny, fill=white, inner sep=0.8pt] at (1.0000,4.1124,2.1124) {$0.04$};
\filldraw[diagcol, fill=diagcol!20, line width=0.6pt] (1.0000,4.0000,2.0000) -- (1.0000,4.3324,2.3324) -- (0.0000,4.0068,2.0068) -- (0.0000,4.0000,2.0000) -- cycle;
\draw[black, line width=1.2pt] (1.0000,4.0000,2.0000) -- (0.0000,4.0000,2.0000);
\node[font=\tiny, fill=white, inner sep=0.8pt] at (1.0000,4.3324,2.3324) {$0.12$};
\node[font=\tiny, fill=white, inner sep=0.8pt] at (0.0000,4.0068,2.0068) {$0.00$};
\filldraw[diagcol, fill=diagcol!20, line width=0.6pt] (0.0000,4.0000,2.0000) -- (0.0000,4.0000,2.1984) -- (0.0000,0.0000,2.1984) -- (0.0000,0.0000,2.0000) -- cycle;
\draw[black, line width=1.2pt] (0.0000,4.0000,2.0000) -- (0.0000,0.0000,2.0000);
\node[font=\tiny, fill=white, inner sep=0.8pt] at (0.0000,4.0000,2.1984) {$0.05$};
\node[font=\tiny, fill=white, inner sep=0.8pt] at (0.0000,0.0000,2.1984) {$0.05$};
\fill (0.0000,0.0000,0.0000) circle (1.3pt);
\node[font=\scriptsize, above left=0pt] at (0.0000,0.0000,0.0000) {0};
\fill (2.0000,0.0000,0.0000) circle (1.3pt);
\node[font=\scriptsize, above left=0pt] at (2.0000,0.0000,0.0000) {1};
\fill (2.0000,4.0000,0.0000) circle (1.3pt);
\node[font=\scriptsize, above left=0pt] at (2.0000,4.0000,0.0000) {2};
\fill (0.0000,4.0000,0.0000) circle (1.3pt);
\node[font=\scriptsize, above left=0pt] at (0.0000,4.0000,0.0000) {3};
\fill (0.0000,0.0000,2.0000) circle (1.3pt);
\node[font=\scriptsize, above left=0pt] at (0.0000,0.0000,2.0000) {4};
\fill (2.0000,0.0000,2.0000) circle (1.3pt);
\node[font=\scriptsize, above left=0pt] at (2.0000,0.0000,2.0000) {5};
\fill (2.0000,4.0000,2.0000) circle (1.3pt);
\node[font=\scriptsize, above left=0pt] at (2.0000,4.0000,2.0000) {6};
\fill (0.0000,4.0000,2.0000) circle (1.3pt);
\node[font=\scriptsize, above left=0pt] at (0.0000,4.0000,2.0000) {7};
\fill (1.0000,0.0000,2.0000) circle (1.3pt);
\node[font=\scriptsize, above left=0pt] at (1.0000,0.0000,2.0000) {8};
\fill (1.0000,4.0000,2.0000) circle (1.3pt);
\node[font=\scriptsize, above left=0pt] at (1.0000,4.0000,2.0000) {9};
\end{tikzpicture}
\caption{Diagrama de $M_x$ [kN\,m], dibujado hacia $y_L$ positivo para valores positivos.}
\label{fig:viga3d-diag-M_x}
\end{figure}

\begin{figure}[!tbp]
\centering
\begin{tikzpicture}[x={(1.5056cm,-0.4346cm)}, y={(0.8693cm,0.7528cm)}, z={(0.0000cm,1.5056cm)}]
\filldraw[diagcol, fill=diagcol!20, line width=0.6pt] (0.0000,0.0000,0.0000) -- (0.3059,0.1530,-0.1530) -- (0.6901,-0.1549,2.1549) -- (1.0000,0.0000,2.0000) -- cycle;
\draw[black, line width=1.2pt] (0.0000,0.0000,0.0000) -- (1.0000,0.0000,2.0000);
\node[font=\tiny, fill=white, inner sep=0.8pt] at (0.3059,0.1530,-0.1530) {$0.10$};
\node[font=\tiny, fill=white, inner sep=0.8pt] at (0.6901,-0.1549,2.1549) {$-0.10$};
\filldraw[diagcol, fill=diagcol!20, line width=0.6pt] (2.0000,0.0000,0.0000) -- (2.2370,0.0000,0.0000) -- (1.7799,0.0000,2.0000) -- (2.0000,0.0000,2.0000) -- cycle;
\draw[black, line width=1.2pt] (2.0000,0.0000,0.0000) -- (2.0000,0.0000,2.0000);
\node[font=\tiny, fill=white, inner sep=0.8pt] at (2.2370,0.0000,0.0000) {$0.06$};
\node[font=\tiny, fill=white, inner sep=0.8pt] at (1.7799,0.0000,2.0000) {$-0.06$};
\filldraw[diagcol, fill=diagcol!20, line width=0.6pt] (2.0000,4.0000,0.0000) -- (2.2370,4.0000,0.0000) -- (1.7799,4.0000,2.0000) -- (2.0000,4.0000,2.0000) -- cycle;
\draw[black, line width=1.2pt] (2.0000,4.0000,0.0000) -- (2.0000,4.0000,2.0000);
\node[font=\tiny, fill=white, inner sep=0.8pt] at (2.2370,4.0000,0.0000) {$0.06$};
\node[font=\tiny, fill=white, inner sep=0.8pt] at (1.7799,4.0000,2.0000) {$-0.06$};
\filldraw[diagcol, fill=diagcol!20, line width=0.6pt] (0.0000,4.0000,0.0000) -- (0.3708,4.1854,-0.1854) -- (0.5473,3.7736,2.2264) -- (1.0000,4.0000,2.0000) -- cycle;
\draw[black, line width=1.2pt] (0.0000,4.0000,0.0000) -- (1.0000,4.0000,2.0000);
\node[font=\tiny, fill=white, inner sep=0.8pt] at (0.3708,4.1854,-0.1854) {$0.12$};
\node[font=\tiny, fill=white, inner sep=0.8pt] at (0.5473,3.7736,2.2264) {$-0.14$};
\filldraw[diagcol, fill=diagcol!20, line width=0.6pt] (0.0000,0.0000,2.0000) -- (0.0000,0.0068,1.9932) -- (1.0000,0.3324,1.6676) -- (1.0000,0.0000,2.0000) -- cycle;
\draw[black, line width=1.2pt] (0.0000,0.0000,2.0000) -- (1.0000,0.0000,2.0000);
\node[font=\tiny, fill=white, inner sep=0.8pt] at (0.0000,0.0068,1.9932) {$0.00$};
\node[font=\tiny, fill=white, inner sep=0.8pt] at (1.0000,0.3324,1.6676) {$0.12$};
\filldraw[diagcol, fill=diagcol!20, line width=0.6pt] (1.0000,0.0000,2.0000) -- (1.0000,0.1124,1.8876) -- (2.0000,0.1168,1.8832) -- (2.0000,0.0000,2.0000) -- cycle;
\draw[black, line width=1.2pt] (1.0000,0.0000,2.0000) -- (2.0000,0.0000,2.0000);
\node[font=\tiny, fill=white, inner sep=0.8pt] at (1.0000,0.1124,1.8876) {$0.04$};
\node[font=\tiny, fill=white, inner sep=0.8pt] at (2.0000,0.1168,1.8832) {$0.04$};
\filldraw[diagcol, fill=diagcol!20, line width=0.6pt] (2.0000,0.0000,2.0000) -- (1.9924,0.0000,2.0000) -- (1.9924,4.0000,2.0000) -- (2.0000,4.0000,2.0000) -- cycle;
\draw[black, line width=1.2pt] (2.0000,0.0000,2.0000) -- (2.0000,4.0000,2.0000);
\node[font=\tiny, fill=white, inner sep=0.8pt] at (1.9924,0.0000,2.0000) {$0.00$};
\node[font=\tiny, fill=white, inner sep=0.8pt] at (1.9924,4.0000,2.0000) {$0.00$};
\filldraw[diagcol, fill=diagcol!20, line width=0.6pt] (2.0000,4.0000,2.0000) -- (2.0000,4.1033,1.8967) -- (1.0000,4.1327,1.8673) -- (1.0000,4.0000,2.0000) -- cycle;
\draw[black, line width=1.2pt] (2.0000,4.0000,2.0000) -- (1.0000,4.0000,2.0000);
\node[font=\tiny, fill=white, inner sep=0.8pt] at (2.0000,4.1033,1.8967) {$-0.04$};
\node[font=\tiny, fill=white, inner sep=0.8pt] at (1.0000,4.1327,1.8673) {$-0.05$};
\filldraw[diagcol, fill=diagcol!20, line width=0.6pt] (1.0000,4.0000,2.0000) -- (1.0000,4.4243,1.5757) -- (0.0000,3.9932,2.0068) -- (0.0000,4.0000,2.0000) -- cycle;
\draw[black, line width=1.2pt] (1.0000,4.0000,2.0000) -- (0.0000,4.0000,2.0000);
\node[font=\tiny, fill=white, inner sep=0.8pt] at (1.0000,4.4243,1.5757) {$-0.15$};
\node[font=\tiny, fill=white, inner sep=0.8pt] at (0.0000,3.9932,2.0068) {$0.00$};
\filldraw[diagcol, fill=diagcol!20, line width=0.6pt] (0.0000,4.0000,2.0000) -- (0.0136,4.0000,2.0000) -- (0.0136,0.0000,2.0000) -- (0.0000,0.0000,2.0000) -- cycle;
\draw[black, line width=1.2pt] (0.0000,4.0000,2.0000) -- (0.0000,0.0000,2.0000);
\node[font=\tiny, fill=white, inner sep=0.8pt] at (0.0136,4.0000,2.0000) {$0.00$};
\node[font=\tiny, fill=white, inner sep=0.8pt] at (0.0136,0.0000,2.0000) {$0.00$};
\fill (0.0000,0.0000,0.0000) circle (1.3pt);
\node[font=\scriptsize, above left=0pt] at (0.0000,0.0000,0.0000) {0};
\fill (2.0000,0.0000,0.0000) circle (1.3pt);
\node[font=\scriptsize, above left=0pt] at (2.0000,0.0000,0.0000) {1};
\fill (2.0000,4.0000,0.0000) circle (1.3pt);
\node[font=\scriptsize, above left=0pt] at (2.0000,4.0000,0.0000) {2};
\fill (0.0000,4.0000,0.0000) circle (1.3pt);
\node[font=\scriptsize, above left=0pt] at (0.0000,4.0000,0.0000) {3};
\fill (0.0000,0.0000,2.0000) circle (1.3pt);
\node[font=\scriptsize, above left=0pt] at (0.0000,0.0000,2.0000) {4};
\fill (2.0000,0.0000,2.0000) circle (1.3pt);
\node[font=\scriptsize, above left=0pt] at (2.0000,0.0000,2.0000) {5};
\fill (2.0000,4.0000,2.0000) circle (1.3pt);
\node[font=\scriptsize, above left=0pt] at (2.0000,4.0000,2.0000) {6};
\fill (0.0000,4.0000,2.0000) circle (1.3pt);
\node[font=\scriptsize, above left=0pt] at (0.0000,4.0000,2.0000) {7};
\fill (1.0000,0.0000,2.0000) circle (1.3pt);
\node[font=\scriptsize, above left=0pt] at (1.0000,0.0000,2.0000) {8};
\fill (1.0000,4.0000,2.0000) circle (1.3pt);
\node[font=\scriptsize, above left=0pt] at (1.0000,4.0000,2.0000) {9};
\end{tikzpicture}
\caption{Diagrama de $M_y$ [kN\,m], dibujado hacia $x_L$ positivo para valores positivos.}
\label{fig:viga3d-diag-M_y}
\end{figure}

\begin{figure}[!tbp]
\centering
\begin{tikzpicture}[x={(1.5293cm,-0.4415cm)}, y={(0.8829cm,0.7646cm)}, z={(0.0000cm,1.5293cm)}]
\filldraw[diagcol, fill=diagcol!20, line width=0.6pt] (0.0000,0.0000,0.0000) -- (-0.0669,-0.0334,0.0334) -- (0.9331,-0.0334,2.0334) -- (1.0000,0.0000,2.0000) -- cycle;
\draw[black, line width=1.2pt] (0.0000,0.0000,0.0000) -- (1.0000,0.0000,2.0000);
\node[font=\tiny, fill=white, inner sep=0.8pt] at (-0.0669,-0.0334,0.0334) {$-0.01$};
\node[font=\tiny, fill=white, inner sep=0.8pt] at (0.9331,-0.0334,2.0334) {$-0.01$};
\filldraw[diagcol, fill=diagcol!20, line width=0.6pt] (2.0000,0.0000,0.0000) -- (1.9822,0.0000,0.0000) -- (1.9822,0.0000,2.0000) -- (2.0000,0.0000,2.0000) -- cycle;
\draw[black, line width=1.2pt] (2.0000,0.0000,0.0000) -- (2.0000,0.0000,2.0000);
\node[font=\tiny, fill=white, inner sep=0.8pt] at (1.9822,0.0000,0.0000) {$-0.00$};
\node[font=\tiny, fill=white, inner sep=0.8pt] at (1.9822,0.0000,2.0000) {$-0.00$};
\filldraw[diagcol, fill=diagcol!20, line width=0.6pt] (2.0000,4.0000,0.0000) -- (2.0178,4.0000,0.0000) -- (2.0178,4.0000,2.0000) -- (2.0000,4.0000,2.0000) -- cycle;
\draw[black, line width=1.2pt] (2.0000,4.0000,0.0000) -- (2.0000,4.0000,2.0000);
\node[font=\tiny, fill=white, inner sep=0.8pt] at (2.0178,4.0000,0.0000) {$0.00$};
\node[font=\tiny, fill=white, inner sep=0.8pt] at (2.0178,4.0000,2.0000) {$0.00$};
\filldraw[diagcol, fill=diagcol!20, line width=0.6pt] (0.0000,4.0000,0.0000) -- (0.0669,4.0334,-0.0334) -- (1.0669,4.0334,1.9666) -- (1.0000,4.0000,2.0000) -- cycle;
\draw[black, line width=1.2pt] (0.0000,4.0000,0.0000) -- (1.0000,4.0000,2.0000);
\node[font=\tiny, fill=white, inner sep=0.8pt] at (0.0669,4.0334,-0.0334) {$0.01$};
\node[font=\tiny, fill=white, inner sep=0.8pt] at (1.0669,4.0334,1.9666) {$0.01$};
\filldraw[diagcol, fill=diagcol!20, line width=0.6pt] (0.0000,0.0000,2.0000) -- (0.0000,0.4243,1.5757) -- (1.0000,0.4243,1.5757) -- (1.0000,0.0000,2.0000) -- cycle;
\draw[black, line width=1.2pt] (0.0000,0.0000,2.0000) -- (1.0000,0.0000,2.0000);
\node[font=\tiny, fill=white, inner sep=0.8pt] at (0.0000,0.4243,1.5757) {$0.05$};
\node[font=\tiny, fill=white, inner sep=0.8pt] at (1.0000,0.4243,1.5757) {$0.05$};
\filldraw[diagcol, fill=diagcol!20, line width=0.6pt] (1.0000,0.0000,2.0000) -- (1.0000,-0.0115,2.0115) -- (2.0000,-0.0115,2.0115) -- (2.0000,0.0000,2.0000) -- cycle;
\draw[black, line width=1.2pt] (1.0000,0.0000,2.0000) -- (2.0000,0.0000,2.0000);
\node[font=\tiny, fill=white, inner sep=0.8pt] at (1.0000,-0.0115,2.0115) {$-0.00$};
\node[font=\tiny, fill=white, inner sep=0.8pt] at (2.0000,-0.0115,2.0115) {$-0.00$};
\filldraw[diagcol, fill=diagcol!20, line width=0.6pt] (2.0000,0.0000,2.0000) -- (2.0000,0.0000,2.0000) -- (2.0000,4.0000,2.0000) -- (2.0000,4.0000,2.0000) -- cycle;
\draw[black, line width=1.2pt] (2.0000,0.0000,2.0000) -- (2.0000,4.0000,2.0000);
\filldraw[diagcol, fill=diagcol!20, line width=0.6pt] (2.0000,4.0000,2.0000) -- (2.0000,3.9885,2.0115) -- (1.0000,3.9885,2.0115) -- (1.0000,4.0000,2.0000) -- cycle;
\draw[black, line width=1.2pt] (2.0000,4.0000,2.0000) -- (1.0000,4.0000,2.0000);
\node[font=\tiny, fill=white, inner sep=0.8pt] at (2.0000,3.9885,2.0115) {$0.00$};
\node[font=\tiny, fill=white, inner sep=0.8pt] at (1.0000,3.9885,2.0115) {$0.00$};
\filldraw[diagcol, fill=diagcol!20, line width=0.6pt] (1.0000,4.0000,2.0000) -- (1.0000,4.4243,1.5757) -- (0.0000,4.4243,1.5757) -- (0.0000,4.0000,2.0000) -- cycle;
\draw[black, line width=1.2pt] (1.0000,4.0000,2.0000) -- (0.0000,4.0000,2.0000);
\node[font=\tiny, fill=white, inner sep=0.8pt] at (1.0000,4.4243,1.5757) {$-0.05$};
\node[font=\tiny, fill=white, inner sep=0.8pt] at (0.0000,4.4243,1.5757) {$-0.05$};
\filldraw[diagcol, fill=diagcol!20, line width=0.6pt] (0.0000,4.0000,2.0000) -- (0.0000,4.0000,2.0000) -- (0.0000,0.0000,2.0000) -- (0.0000,0.0000,2.0000) -- cycle;
\draw[black, line width=1.2pt] (0.0000,4.0000,2.0000) -- (0.0000,0.0000,2.0000);
\fill (0.0000,0.0000,0.0000) circle (1.3pt);
\node[font=\scriptsize, above left=0pt] at (0.0000,0.0000,0.0000) {0};
\fill (2.0000,0.0000,0.0000) circle (1.3pt);
\node[font=\scriptsize, above left=0pt] at (2.0000,0.0000,0.0000) {1};
\fill (2.0000,4.0000,0.0000) circle (1.3pt);
\node[font=\scriptsize, above left=0pt] at (2.0000,4.0000,0.0000) {2};
\fill (0.0000,4.0000,0.0000) circle (1.3pt);
\node[font=\scriptsize, above left=0pt] at (0.0000,4.0000,0.0000) {3};
\fill (0.0000,0.0000,2.0000) circle (1.3pt);
\node[font=\scriptsize, above left=0pt] at (0.0000,0.0000,2.0000) {4};
\fill (2.0000,0.0000,2.0000) circle (1.3pt);
\node[font=\scriptsize, above left=0pt] at (2.0000,0.0000,2.0000) {5};
\fill (2.0000,4.0000,2.0000) circle (1.3pt);
\node[font=\scriptsize, above left=0pt] at (2.0000,4.0000,2.0000) {6};
\fill (0.0000,4.0000,2.0000) circle (1.3pt);
\node[font=\scriptsize, above left=0pt] at (0.0000,4.0000,2.0000) {7};
\fill (1.0000,0.0000,2.0000) circle (1.3pt);
\node[font=\scriptsize, above left=0pt] at (1.0000,0.0000,2.0000) {8};
\fill (1.0000,4.0000,2.0000) circle (1.3pt);
\node[font=\scriptsize, above left=0pt] at (1.0000,4.0000,2.0000) {9};
\end{tikzpicture}
\caption{Diagrama de $T$ [kN\,m], dibujado hacia $x_L$ positivo para valores positivos.}
\label{fig:viga3d-diag-T}
\end{figure}

\section{Deformaciones y tensiones}

\begin{center}
\begin{minipage}{\linewidth}
\centering
\captionof{table}{Deformaciones generalizadas de la sección: $\varepsilon=N/EA$, $\kappa=M/EI$, $\theta'=T/GJ$.}\label{tab:viga3d-deformaciones}
\begin{tabular}{@{}cccccc@{}}
\toprule
Elem. & Nodo & $\varepsilon$ [$\mu\varepsilon$] & $\kappa_x$ [$10^{-3}$/m] & $\kappa_y$ [$10^{-3}$/m] & $\theta'$ [mrad/m] \\
\midrule
0 & 0 & $-2.07$ & $4.2537$ & $5.8142$ & $-0.5465$ \\
 & 8 & $-1.32$ & $-6.2764$ & $-5.8884$ & $-0.5465$ \\
1 & 1 & $-1.65$ & $2.4705$ & $3.6769$ & $-0.1187$ \\
 & 5 & $-0.90$ & $-4.9419$ & $-3.4151$ & $-0.1187$ \\
2 & 2 & $-1.65$ & $-2.4705$ & $3.6769$ & $0.1187$ \\
 & 6 & $-0.90$ & $4.9419$ & $-3.4151$ & $0.1187$ \\
3 & 3 & $-2.07$ & $1.4978$ & $7.0467$ & $0.5465$ \\
 & 9 & $-1.32$ & $-0.2047$ & $-8.6038$ & $0.5465$ \\
4 & 4 & $0.00$ & $-0.1494$ & $0.1494$ & $4.0025$ \\
 & 8 & $0.00$ & $9.3100$ & $7.2931$ & $4.0025$ \\
5 & 8 & $-0.28$ & $2.9114$ & $2.4661$ & $-0.1088$ \\
 & 5 & $-0.28$ & $2.2664$ & $2.5633$ & $-0.1088$ \\
6 & 5 & $-0.34$ & $5.0257$ & $0.1187$ & $0.0000$ \\
 & 6 & $-0.34$ & $5.0257$ & $0.1187$ & $0.0000$ \\
7 & 6 & $-0.28$ & $2.5633$ & $-2.2664$ & $0.1088$ \\
 & 9 & $-0.28$ & $2.4661$ & $-2.9114$ & $0.1088$ \\
8 & 9 & $0.00$ & $7.2931$ & $-9.3100$ & $-4.0025$ \\
 & 7 & $0.00$ & $0.1494$ & $0.1494$ & $-4.0025$ \\
9 & 7 & $-0.13$ & $3.0788$ & $0.2113$ & $0.0000$ \\
 & 4 & $-0.13$ & $3.0788$ & $0.2113$ & $0.0000$ \\
\bottomrule
\end{tabular}
\end{minipage}
\end{center}

\begin{center}
\begin{minipage}{\linewidth}
\centering
\captionof{table}{Tensiones en los extremos: $\sigma_N=N/A$, $\sigma_{\max}=|N|/A+|M_x|c_y/I_x+|M_y|c_x/I_y$, $\tau_T=|T|r/J$, $\sigma_{\mathrm{VM}}=\sqrt{\sigma_{\max}^2+3\tau_T^2}$.}\label{tab:viga3d-tensiones}
\begin{tabular}{@{}cccccc@{}}
\toprule
Elem. & Nodo & $\sigma_N$ [MPa] & $\sigma_{\max}$ [MPa] & $\tau_T$ [MPa] & $\sigma_{\mathrm{VM}}$ [MPa] \\
\midrule
0 & 0 & $-0.43$ & $37.43$ & $0.77$ & $37.45$ \\
 & 8 & $-0.27$ & $44.98$ & $0.77$ & $45.00$ \\
1 & 1 & $-0.34$ & $22.93$ & $0.17$ & $22.94$ \\
 & 5 & $-0.19$ & $30.90$ & $0.17$ & $30.90$ \\
2 & 2 & $-0.34$ & $22.93$ & $0.17$ & $22.94$ \\
 & 6 & $-0.19$ & $30.90$ & $0.17$ & $30.90$ \\
3 & 3 & $-0.43$ & $31.83$ & $0.77$ & $31.86$ \\
 & 9 & $-0.27$ & $32.65$ & $0.77$ & $32.67$ \\
4 & 4 & $0.00$ & $1.10$ & $5.66$ & $9.86$ \\
 & 8 & $0.00$ & $61.02$ & $5.66$ & $61.80$ \\
5 & 8 & $-0.06$ & $19.82$ & $0.15$ & $19.82$ \\
 & 5 & $-0.06$ & $17.81$ & $0.15$ & $17.81$ \\
6 & 5 & $-0.07$ & $18.98$ & $0.00$ & $18.98$ \\
 & 6 & $-0.07$ & $18.98$ & $0.00$ & $18.98$ \\
7 & 6 & $-0.06$ & $17.81$ & $0.15$ & $17.81$ \\
 & 9 & $-0.06$ & $19.82$ & $0.15$ & $19.82$ \\
8 & 9 & $0.00$ & $61.02$ & $5.66$ & $61.80$ \\
 & 7 & $0.00$ & $1.10$ & $5.66$ & $9.86$ \\
9 & 7 & $-0.03$ & $12.12$ & $0.00$ & $12.12$ \\
 & 4 & $-0.03$ & $12.12$ & $0.00$ & $12.12$ \\
\bottomrule
\end{tabular}
\end{minipage}
\end{center}

El elemento más solicitado es el \textbf{4} (nodos 4--8), con $\sigma_{\mathrm{VM}}=61.80$~MPa.

\section{Reacciones y equilibrio}

\begin{center}
\begin{minipage}{\linewidth}
\centering
\captionof{table}{Reacciones en los apoyos (ejes globales).}\label{tab:viga3d-reacciones}
\begin{tabular}{@{}ccccccc@{}}
\toprule
Nodo & $R_x$ [kN] & $R_y$ [kN] & $R_z$ [kN] & $M_x$ [kN\,m] & $M_y$ [kN\,m] & $M_z$ [kN\,m] \\
\midrule
0 & $0.058$ & $0.035$ & $0.448$ & $-0.019$ & $-0.115$ & $0.017$ \\
1 & $-0.058$ & $0.061$ & $0.340$ & $-0.040$ & $-0.060$ & $0.001$ \\
2 & $-0.058$ & $-0.061$ & $0.340$ & $0.040$ & $-0.060$ & $-0.001$ \\
3 & $0.058$ & $-0.035$ & $0.448$ & $0.019$ & $-0.115$ & $-0.017$ \\
\bottomrule
\end{tabular}
\end{minipage}
\end{center}

\begin{equation*}
\textstyle\sum F=(-0.000,\;0.000,\;-1.576)~\text{kN},\qquad \sum R=(-0.000,\;0.000,\;1.576)~\text{kN}
\end{equation*}

El residuo máximo es $4.5\times 10^{-14}$~kN, de modo que se cumple el equilibrio global. La energía de deformación es $\tfrac12\mathbf F^T\mathbf u=2.6469$~J.

\end{document}